\documentclass[10pt,a4paper]{amsart}
\usepackage[margin=19mm]{geometry}
\usepackage{amsmath,amssymb,mathtools}
\usepackage[scr=rsfs,cal=euler]{mathalfa}
\usepackage{xfrac}
\usepackage[unicode,hidelinks,bookmarks=false]{hyperref}
\newcommand{\ie}{i.e.\ }
\newcommand{\cf}{cf.\ }
\newcommand{\resp}{respectively\ }
\newcommand{\Subsection}{\subsection}
\providecommand{\nobreakdash}{\nobreak\hbox{-}}
\setlength{\emergencystretch}{2em}
\multlinegap0pt
\usepackage{amssymb}
\usepackage{mathtools}
\usepackage{bm}
\usepackage{tikz}
\usepackage[shortlabels]{enumitem}
\newtheorem{lemma}{Lemma}
\title{\vspace{-2em} {\Huge\bfseries A note on coextensivity of bounded hoops}\\[1em] \rule{0.75\textwidth}{0.6pt}}
\author{Michael Hoefnagel, Giuseppe Metere, Danielle Kleyn}
\date{Source: arXiv:2608.30476v2 (CC BY 4.0)}
\begin{document}
\maketitle

\noindent\emph{Source and licence.} \vspace{-2em} {\Huge\bfseries A note on coextensivity of bounded hoops}\\[1em] \rule{0.75\textwidth}{0.6pt}. Version arXiv:2608.30476v2 (CC BY 4.0), distributed by arXiv under the Creative Commons Attribution 4.0 International licence, http://creativecommons.org/licenses/by/4.0/, as recorded in the arXiv licence metadata for this exact version and shown on the abstract page https://arxiv.org/abs/2608.30476v2. Full licence notice: \texttt{licenses/arxiv-2608.30476-CC-BY-4.0.txt}.\par\smallskip

\noindent\textit{Editorial scope.} Counted unit: the article's lemma central (source0.tex lines 183--189). The article's complete original proof of it is delivered with it (source0.tex lines 190--206, one \texttt{proof} environment of 17 lines). No mathematics is written, repaired or re-derived: the statement and the proof are the article's own lines, in the article's own notation, and no retained line is substituted or re-flowed.  Retained uncounted as the source-local context the proof needs: lines 154--156.  Nothing was omitted from any retained span, so the package carries no omission record.  No retained line contains a citation, so the wrapper carries no bibliography and no bibliography entry.  No retained line contains a figure environment, an image-inclusion command or drawing code, so no image is delivered and the package declares no asset dependency.  Editorial formatting only, in addition: an AMS \texttt{amsart} preamble that restates the article's own declarations and macros used by the retained lines, namely the environments \texttt{lemma} on separate counters, together with the standard packages those lines need.  The retained environments print in this document as Lemma 1 in order of appearance, whereas the article prints the same items with the numbering of its own full document.  The complete native source of the article as distributed in the arXiv e-print for this exact version is bundled unchanged as \texttt{sources/208944/source0.tex}.
\begin{lemma} \label{lemma: idempotent}
Let $e$ be any idempotent in a hoop $X$, then $e\wedge x = e\cdot x$ and $e \wedge (-)$ is left adjoint to $e \to (-)$.
\end{lemma}

\begin{lemma} \label{lemma: weakly complementary central}
Let $X$ be any bounded hoop and $e_1, e_2$ any two weakly complementary central elements. Then the induced map 
\[
X \to [e_1 \to X] \times [e_2 \to X]
\]
is surjective. It is injective if and only if $e_1$ and $e_2$ admit a join and that join is $1$.
\end{lemma}

\begin{proof}
By Lemma~\ref{lemma: idempotent}, the map $x \mapsto e\to x$ preserves finite meets, since it is a right adjoint.
For surjectivity, note that for any pair $(e_1 \to x, e_2 \to y)$ we may define $z = (e_1 \to x) \wedge (e_2 \to y)$ so that 
\begin{align*}
    e_1 \to z &= e_1 \to \big ((e_1 \to x) \wedge (e_2 \to y)\big) \\ 
    &= (e_1 \to (e_1 \to x))\wedge (e_1 \to (e_2 \to y)) \\
    &= (e_1 \to x) \wedge (e_1\cdot e_2 \to y) \\
    &= (e_1 \to x) \wedge (0 \to y) = e_1 \to x
\end{align*}
and similarly $e_2 \to z = e_2 \to y$, hence $(e_1 \to z, e_2 \to z) = (e_1 \to x, e_2 \to y)$. 

For injectivity, if $e_1 \vee e_2 = 1$ then $(e_1 \to x, e_2 \to x) = (e_1 \to y, e_2 \to y)$ implies that 
\begin{align*}
(e_i \to x) \to (e_i \to y) = 1 \implies e_i \to (x \to y) = 1 \implies e_i \leqslant x \to y
\end{align*}
so that $1 = e_1 \vee e_2 \leqslant x \to y$, which implies $x \leqslant y$ and similarly $y \leqslant x$. On the other hand, if $x \mapsto (e_1 \to x, e_2 \to x)$ is injective, and hence an isomorphism, then note that under this map $e_1$ is sent to $(1, e_1)$ and $e_2$ is sent to $(e_2, 1)$, whose join is $(1,1)$ and hence $e_1 \vee e_2 = 1$. 
\end{proof}

\end{document}
